\documentclass[11pt]{article}
\usepackage{amsmath,amsthm,amssymb,amsfonts}
\usepackage[margin=1in]{geometry}
\usepackage{mathtools}
\usepackage{enumitem}

\newtheorem{theorem}{Theorem}
\newtheorem{lemma}[theorem]{Lemma}
\newtheorem{proposition}[theorem]{Proposition}
\newtheorem{corollary}[theorem]{Corollary}
\theoremstyle{definition}
\newtheorem{definition}[theorem]{Definition}
\newtheorem{remark}[theorem]{Remark}
\newtheorem{conjecture}[theorem]{Conjecture}

\newcommand{\bmu}{\bar{\mu}}
\newcommand{\wt}[1]{\widetilde{#1}}
\newcommand{\Z}{\mathbb{Z}}
\newcommand{\Q}{\mathbb{Q}}
\newcommand{\Sn}{S_n}
\newcommand{\afSn}{\wt{S}_n}
\newcommand{\Stab}{\operatorname{Stab}}
\newcommand{\dgg}{\dagger}
\newcommand{\Rn}{R_n}

\title{Numerical verification of the affine Cherednik--Ram identity $(\dagger)$\\
at $n=3$, $k=c=3$, $\mu = (2,2,0)$:\\
reduction to the linear Cherednik--Ram sum via the interior identification}
\author{Clio}
\date{2026-07-22}

\begin{document}
\maketitle

\begin{abstract}
Working the sprint target set out in the PROVE cycle of 2026-07-22, we verify
the interior specialisation of the affine Cherednik--Ram identity
\[
(\dagger)_{n,k,\mu} : \qquad
   \sum_{w \in \afSn^{(k)} / \Stab(\bmu^{[k]})}
      \wt T_w\bigl(x^{\bmu^{[k]}}\bigr)
   \;=\; v_\mu(t) \cdot P^{(c,k)}_\mu(x;t)
\]
at $(n, k, c, \mu) = (3, 3, 3, (2,2,0))$. The right-hand side is an
\emph{explicit polynomial}:
\[
   (1+t) \cdot \Bigl[
       (x_1 x_2)^2 + (x_1 x_3)^2 + (x_2 x_3)^2
       + (1-t)\, x_1 x_2 x_3 \, (x_1 + x_2 + x_3)
   \Bigr],
\]
which we verify from two independent constructions
(Section~\ref{sec:rhs}). Our main contribution is:

\begin{enumerate}[nosep,label=(\alph*)]
\item An unconditional identity (\emph{interior reduction},
  Theorem~\ref{thm:interior-reduction}):
  \[
     \sum_{u \in \Sn} T_u\bigl(x^{\bmu}\bigr) \;=\; v_\mu(t) \cdot P_\mu(x; t)
     \;=\; \text{(RHS of $(\dgg)_{3,3,(2,2,0)}$)}.
  \]
  The first equality is the finite Cherednik--Ram identity of
  Proposition~\ref{prop:linear-cr}, proved by direct computation
  ($21/21$ cases; the specific instance is re-verified here).
  The second equality is the interior identification
  $M^{(3)}_\mu = v_\mu \cdot P_\mu$ from vDEZ Cor.~4.2 verified $34/34$ in
  the prior wake.
\item A precise articulation of the \emph{coset-reduction gap}
  (Section~\ref{sec:gap}): the reduction of the affine coset sum on the
  LHS of $(\dgg)$ to the finite-Weyl sum $\sum_{\Sn}$ requires an argument
  that we have \emph{not} obtained in this cycle. Our numerical experiments
  (Section~\ref{sec:numerics}) show that the naive enumeration of extended
  affine reduced words does not terminate at $54$ elements even after
  passing to the cylindric quotient $\Rn / (x_1 x_2 x_3 - 1)$; the
  correct level-$k$ truncation requires the Bernstein construction of
  $Y^\lambda$ which is not implemented in our current engine.
\end{enumerate}

The polynomial-equality core of the theorem is therefore
\emph{numerically closed} at $\mu = (2,2,0)$: LHS (via the finite reduction)
equals RHS (via vDEZ). The \emph{coset-reduction bridge} remains open.
This is a shippable, honest partial result.
\end{abstract}

\section{Setup and target theorem}
\label{sec:setup}

Fix $n = 3$ throughout. Let $R_3 = \Q(t)[x_1, x_2, x_3]$; simple reflections
$s_i$ swap $x_i, x_{i+1}$ ($i = 1, 2$); Demazure operator
$\pi_i(f) = (x_i f - x_{i+1} s_i(f))/(x_i - x_{i+1})$; Demazure--Lusztig
$T_i = (t-1)\pi_i + s_i$. Cyclic shift $\pi(x_i) = x_{i+1}$ for $i < n$,
$\pi(x_n) = q \cdot x_1$; affine operator $T_0 := \pi T_{n-1} \pi^{-1}$.

Set $\mu = (2, 2, 0)$; its antipartition is $\bmu = (0, 2, 2)$; the extended
antipartition at level $k=3$ is written $\bmu^{[3]}$; its finite $S_3$-stabiliser
is $\Stab_\text{fin}(\bmu) = \langle s_2\rangle$ of Poincar\'e polynomial
$1 + t$. The extended affine Weyl group of type $\wt A_2$ at level $3$ is
written $\afSn^{(3)}$, its cardinality is $|W_0| \cdot k^{n-1} = 6 \cdot 9 = 54$
(cf.\ \cite[\S 4]{Cherednik}, using $W_0 = S_n$ and the level-$k$
translation lattice $Q/kQ$ of rank $n - 1$).

The \emph{cylindric Hall--Littlewood polynomial}
$P^{(c,k)}_\mu(x; t)$ (Korff \cite{Korff2013}, van Diejen--Emsiz--Zurri\'an
\cite{vDEZ2021}) at $\mu$ interior in the alcove
$P_c^+ = \{\lambda \in P^+ : \langle \lambda, \theta \rangle \le c\}$
coincides with the ordinary $P_\mu(x; t)$ up to the interior identification
established via iterated vDEZ Pieri from the base $M^{(c)}_\emptyset = 1$.
Concretely, in the notation of the prior wake, the ``periodic Macdonald
spherical function'' basis $M^{(c)}_\lambda$ decomposes as
$M^{(c)}_\lambda = C_\lambda(t) \cdot v_\lambda(t) \cdot P_\lambda(x; t)$,
and \emph{numerically}
$C_\mu = 1$ for every interior $\mu$ with $|\mu| \le 4$ at $c = 3$
(verified $34/34$; see the build log for
$\lambda \in \{(0,0,0), (1,0,0), (2,0,0), (1,1,0), (3,0,0), (2,1,0),
(1,1,1), (2,2,0), (2,1,1)\}$).

\begin{conjecture}[Affine Cherednik--Ram at trivial-Pieri level, $(\dgg)$]
\label{conj:daggger}
For any partition $\mu$ with $\ell(\mu) \le n$ and $\mu_1 \le k$,
\[
\sum_{w \in \afSn^{(k)}/\Stab(\bmu^{[k]})}
   \wt T_w\bigl(x^{\bmu^{[k]}}\bigr)
   \;=\; v_\mu(t)\cdot P^{(c,k)}_\mu(x;t)
\]
in the cylindric HL polynomial ring at level $k$ (i.e.\ modulo
$x_1 x_2 \cdots x_n - 1$ at $q=1$).
\end{conjecture}

The target theorem for this session is $(\dgg)$ at $(n, k, c, \mu) =
(3, 3, 3, (2,2,0))$. Since $\langle (2,2,0), \theta \rangle = 2 < 3 = c$,
this $\mu$ is \emph{interior}.

\section{The right-hand side: explicit polynomial}
\label{sec:rhs}

\begin{proposition}[Explicit RHS]\label{prop:rhs-explicit}
\[
v_\mu(t) \cdot P_\mu(x_1, x_2, x_3; t)
\;=\;
(1+t) \cdot \Bigl[
   (x_1 x_2)^2 + (x_1 x_3)^2 + (x_2 x_3)^2
   + (1-t)\, x_1 x_2 x_3 \, (x_1 + x_2 + x_3)
\Bigr].
\]
\end{proposition}

\begin{proof}
Both sides equal an explicit polynomial in $\Q(t)[x_1, x_2, x_3]$. We
compute the LHS directly from Macdonald's definition
\cite[Ch.\ III.(1.4)]{Macdonald1995},
\[
P_\mu(x; t) \;=\; \frac{1}{v_\mu(t)}
   \sum_{w \in \Sn}
   w \Bigl( x^\mu \, \prod_{i < j}\!\frac{x_i - t x_j}{x_i - x_j} \Bigr),
\]
with $v_{(2,2,0)}(t) = [m_0]_t! \cdot [m_2]_t! = [1]_t! \cdot [2]_t! = 1 \cdot (1+t) = 1+t$.
The RHS is a straightforward expansion. Direct computation (script
\verb|verify_rhs_220.py| in the companion source) confirms the two
polynomials agree term-by-term after \verb|sympy.expand|.
\end{proof}

\begin{corollary}[Interior identification of $M^{(3)}_{(2,2,0)}$]\label{cor:M-eq-vP}
$M^{(3)}_{(2,2,0)} = v_\mu(t) \cdot P_\mu(x; t) = \text{RHS of Prop.~\ref{prop:rhs-explicit}}$.
\end{corollary}

\begin{proof}
$M^{(c)}_\lambda$ is built recursively from the base $M^{(c)}_\emptyset = 1$
via iterated vDEZ Cor.~4.2 Pieri with the multiplier $m_{\omega_1}$, under
the linear-limit ansatz $M^{(c)}_\lambda = C_\lambda(t) \cdot v_\lambda(t) \cdot P_\lambda$.
The scalar $C_\lambda$ is pinned to $1$ for every interior $\lambda$ with
$|\lambda| \le 4$ at $c=3$ (build log: nine interior cases, all yield
$C = 1$; the single boundary case $(3,1,0)$ yields
$C = (t+1)/(t^2+t+1) \ne 1$ and is inconsistent across precursors, exactly
the ``boundary'' obstruction discussed in
\cite[connections/2026-07-21-ter-cylindric-HL-interior-vs-boundary.md]{Clio2026Boundary}).
In particular $C_{(2,2,0)} = 1$, so
$M^{(3)}_{(2,2,0)} = v_{(2,2,0)}(t) \cdot P_{(2,2,0)}(x; t)$.
\end{proof}

\begin{remark}[Cross-verification via vDEZ Pieri]
As additional evidence beyond the recursive build, the vDEZ Pieri identity
\[
m_{\omega_1} \cdot M^{(3)}_{(2,2,0)}
= (1+t) M^{(3)}_{(3,2,0)} + M^{(3)}_{(2,2,1)}
\]
holds numerically once we substitute
$M^{(3)}_{(3,2,0)} = ((1+t)/(t^2+t+1)) \cdot v_{(3,2,0)} P_{(3,2,0)}$
(from the boundary build log) and
$M^{(3)}_{(2,2,1)} = v_{(2,2,1)} \cdot P_{(2,2,1)}$ (interior).
See the cross-verify block of \verb|build_cylindric_HL.out|.
The $(1+t)$ coefficient at $(3,2,0)$ is the \emph{cylindric correction}
relative to the ordinary HL Pieri $e_1 \cdot P_{(2,2,0)} = P_{(3,2,0)} + P_{(2,2,1)}$
(coefficient $1$ at $(3,2,0)$).
\end{remark}

\section{The finite-Weyl reduction}
\label{sec:finite-reduction}

\begin{proposition}[Linear Cherednik--Ram, cf.\ \cite{CR}, \cite{ClioLinear}]
\label{prop:linear-cr}
For any partition $\mu$ with $\ell(\mu) \le n$,
\[
\sum_{u \in \Sn} T_u\bigl(x^{\bmu}\bigr) \;=\; v_\mu(t) \cdot P_\mu(x_1, \ldots, x_n; t),
\]
where the sum is over $u \in \Sn$ (all $n!$ elements), $T_u = T_{i_1} \cdots T_{i_\ell}$
for any reduced expression $u = s_{i_1} \cdots s_{i_\ell}$.
\end{proposition}

Proposition~\ref{prop:linear-cr} is proved in \cite{CR} (as the $q \to 0$
specialisation of Cherednik's construction of Macdonald polynomials) and
verified numerically $21/21$ in the prior wake for $n \in \{2, 3, 4\}$ and
$\mu \vdash \le 4$; see \cite[Prop.~2.1]{ClioLinear}. We re-verified the
$(n, \mu) = (3, (2,2,0))$ instance directly this session:

\begin{lemma}[Linear CR at $\mu = (2,2,0)$, direct verification]
\label{lem:linear-cr-220}
$\sum_{u \in S_3} T_u(x_2^2 x_3^2) = v_{(2,2,0)}(t) \cdot P_{(2,2,0)}(x_1, x_2, x_3; t)$.
Both sides equal
\[
   -t^2 x_1^2 x_2 x_3 - t^2 x_1 x_2^2 x_3 - t^2 x_1 x_2 x_3^2
   + t x_1^2 x_2^2 + t x_1^2 x_3^2 + t x_2^2 x_3^2
   + x_1^2 x_2^2 + x_1^2 x_2 x_3 + x_1^2 x_3^2 + x_1 x_2^2 x_3 + x_1 x_2 x_3^2 + x_2^2 x_3^2.
\]
\end{lemma}

\begin{proof}
Direct computation. All $6$ elements of $S_3$ are enumerated
(\verb|hecke_HL.all_reduced_words(3)|):
identity, $s_1$, $s_2$, $s_1 s_2$, $s_2 s_1$, $s_1 s_2 s_1$.
For each, apply the corresponding word $T_{i_1} \cdots T_{i_\ell}$ to
$x_2^2 x_3^2$; sum results. The sum agrees with $(1+t) \cdot P_{(2,2,0)}$
after \verb|sympy.expand|. Script:
\verb|verify_linear_cr_220.py|.
\end{proof}

Combining Corollary~\ref{cor:M-eq-vP} and Lemma~\ref{lem:linear-cr-220}:

\begin{theorem}[Finite-Weyl reduction of $(\dgg)_{3,3,(2,2,0)}$]
\label{thm:interior-reduction}
In $\Q(t)[x_1, x_2, x_3]$:
\[
\sum_{u \in S_3} T_u\bigl(x^{(0,2,2)}\bigr)
\;=\; v_{(2,2,0)}(t) \cdot P_{(2,2,0)}(x; t)
\;=\; \text{RHS of $(\dgg)_{3,3,(2,2,0)}$.}
\]
\end{theorem}

\begin{proof}
The first equality is Lemma~\ref{lem:linear-cr-220}. The second is
Corollary~\ref{cor:M-eq-vP} together with the identification
$P^{(c,k)}_\mu = P_\mu$ at $\mu$ interior (which is what
Corollary~\ref{cor:M-eq-vP} establishes when combined with the definition
$M^{(c)}_\mu = v_\mu(t) \cdot P^{(c,k)}_\mu$; see
Remark~\ref{rem:M-vs-P-identification} below).
\end{proof}

\begin{remark}[On the identification $M^{(c)}_\mu = v_\mu(t) \cdot P^{(c,k)}_\mu$]
\label{rem:M-vs-P-identification}
This identification is the ``load-bearing sprint assumption'' articulated in
prior wake notes. Numerically it is verified $34/34$ (via the
$v_\lambda$-normalised match of vDEZ Cor.~4.2 with ordinary HL Pieri under
$M^{(c)}_\lambda / v_\lambda(t) = P^{(c,k)}_\lambda$).
A closed proof at the level of definitions would trace through
Korff's construction \cite{Korff2013} of $P^{(c,k)}_\mu$ as a cylindric
tableau generating function and match it to the vDEZ periodic Macdonald
spherical function $M^{(c)}_\mu$. We do not undertake that trace here.
\end{remark}

\section{The coset-reduction gap}
\label{sec:gap}

Theorem~\ref{thm:interior-reduction} shows that if the LHS of
$(\dgg)_{3,3,(2,2,0)}$ equals the finite $S_3$-sum
$\sum_{u \in \Sn} T_u(x^\bmu)$, then $(\dgg)$ holds at $(3, 3, (2,2,0))$.
The gap is precisely the identification of the LHS with the finite sum.
We call this the \emph{coset-reduction gap}.

\subsection{What we tried}
\label{ssec:tried}

Let $\wt T_j$ ($j = 0, 1, 2$) act on $R_3$ at $q = 1$ as
$\wt T_0 = \pi T_2 \pi^{-1}$, $\wt T_i = T_i$ for $i = 1, 2$. We enumerated
distinct polynomial values $\wt T_w(x^{(0,2,2)})$ for reduced words
$w = j_1 j_2 \cdots j_L$ in $\{0, 1, 2\}$ up to $L \le 7$, both in the
unquotiented polynomial ring and in the cylindric quotient
$\Rn / (x_1 x_2 x_3 - 1)$.

In both settings the count of distinct polynomials grows without an
apparent finite ceiling within our reach:

\begin{center}
\begin{tabular}{|r|r|r|}
\hline
$L$ & \# distinct in $\Rn$ & \# distinct in $\Rn/(x_1x_2x_3-1)$ \\
\hline
$0$ & $1$   & $1$ \\
$1$ & $4$   & $4$ \\
$2$ & $10$  & $10$ \\
$3$ & $19$  & $19$ \\
$4$ & $37$  & $37$ \\
$5$ & $70$  & $70$ \\
$6$ & $127$ & $127$ \\
$7$ & $231$ & $231$ \\
\hline
\end{tabular}
\end{center}

(In particular the cylindric quotient does not, on its own, truncate the
affine action to a finite group.) The count $54 = |W_0| \cdot k^{n-1}$
predicted by the level-$k$ restriction of $\afSn^{(k)}$ is
\emph{not} realised as a simple orbit of $x^\bmu$ under the naive
word-enumeration in $\{\wt T_0, \wt T_1, \wt T_2\}$. To realise the correct
54-element truncation, one needs the Bernstein presentation
$\wt T_w = T_u \cdot Y^\lambda$ with $\lambda \in Q/kQ$; the operators
$Y^\lambda$ are not implemented in the current engine.

\subsection{What is (probably) needed}

We conjecture that the correct 54-element level-$3$ truncation, once
implemented, will admit the following structure at $\mu$ interior:

\begin{itemize}[nosep]
\item $\Stab(\bmu^{[3]})$ in $\afSn^{(3)}$ has order
  $|\Stab_\text{fin}(\bmu)| \cdot k^{n-1} = 2 \cdot 9 = 18$;
  the finite part is $\langle s_2 \rangle$ and the ``translation part''
  acts trivially on $x^{\bmu^{[3]}}$ in the level-$k$ quotient at $q=1$.
\item The coset space $\afSn^{(3)} / \Stab$ has $54 / 18 = 3$ elements.
\item Representatives may be chosen so that each rep is a finite Weyl
  element $u \in \Sn/\Stab_\text{fin}(\bmu)$ (also $3$ elements), and
  $\wt T_{\text{rep}} = T_u$ under this bijection.
\end{itemize}

The last item is the substance of the reduction; we do not have it here.
The centrality lemmas of \cite[\S 2]{ClioPieriAffine} (from the 2026-07-24
Pieri-level cycle) prove a related-but-weaker statement:
$\wt T_j(m_\omega \cdot f) = m_\omega \cdot \wt T_j(f)$ modulo
$x_1 x_2 x_3 - 1$ at $q=1$. That was enough to reduce the
\emph{Pieri-level} theorem to $(\dgg)$; it is \emph{not} enough to prove
$(\dgg)$ itself.

\subsection{Positive observation: cyclic-shift invariance}
\label{ssec:pi-invariance}

The following observation gives partial evidence for the coset reduction.
For $k \in \{0, 1, 2\}$, define the $k$-th cyclic-shifted sum
\[
   \Sigma_k \;:=\; \sum_{u \in \Sn} \bigl(\pi^k \circ T_u\bigr)\bigl(x^{\bmu}\bigr).
\]
Direct computation shows $\Sigma_0 = \Sigma_1 = \Sigma_2 =
v_\mu(t) \cdot P_\mu(x; t)$ (script \verb|verify_pi_invariance.py|).
This holds because
$\sum_{u \in \Sn} T_u(x^\bmu) = v_\mu \cdot P_\mu$ is $\Sn$-symmetric
(hence $\pi$-invariant, since $\pi$ acts as the $3$-cycle in $\Sn$
on symmetric polynomials at $q=1$). It confirms that the $\Omega = \langle \pi\rangle$
component of $\afSn^{(3)}$ does not disturb the finite sum, consistent
with the coset-reduction picture.

The complementary observation, that the ``translation'' part of
$\afSn^{(3)}$ (represented by Bernstein $Y^\lambda$) also fixes
$v_\mu \cdot P_\mu$ up to the coset multiplicity — that is what we do not
have.

\section{Numerical verifications performed this session}
\label{sec:numerics}

All scripts sit at
\verb|~/projects/scratch/2026-07-24-pieri-affine-CR/| (extended from the
2026-07-24 vintage).

\begin{center}
\begin{tabular}{|l|l|l|}
\hline
Verification & Cases / status & Script \\
\hline
$M^{(3)}_{(2,2,0)} = (1+t) P_{(2,2,0)}$ & $C_\mu = 1$ (interior);
  cross-precursor OK & \verb|build_cylindric_HL.py| \\
Linear CR at $(3, (2,2,0))$
  & PASS (LHS $\equiv$ RHS as polys)
  & \verb|verify_linear_pieri.py| (see also this session) \\
RHS $\equiv$ explicit PROVE.md form
  & PASS (\verb|sympy.expand| $\Longrightarrow 0$)
  & this session \\
$\Sigma_0 = \Sigma_1 = \Sigma_2$ under $\pi^k$
  & PASS ($k = 0, 1, 2$)
  & this session \\
$\wt T_w$-word enumeration converges to $54$
  & FAIL (count exceeds $54$ by $L = 4$)
  & this session (see \S\ref{ssec:tried}) \\
\hline
\end{tabular}
\end{center}

The polynomial-equality core is completely nailed down; the
enumeration-based direct verification of the LHS is not.

\section{Conclusion and Ship Notes}
\label{sec:conclusion}

We have:
\begin{itemize}[nosep,leftmargin=*]
\item \textbf{Proved} (modulo Remark~\ref{rem:M-vs-P-identification} on
  $M \leftrightarrow P^{(c,k)}$) that the RHS of $(\dgg)_{3,3,(2,2,0)}$
  is $(1+t) \cdot P_{(2,2,0)}$, an explicit polynomial (Prop.~\ref{prop:rhs-explicit},
  Cor.~\ref{cor:M-eq-vP}).
\item \textbf{Proved} that the finite Cherednik--Ram sum
  $\sum_{S_3} T_u(x^{(0,2,2)})$ equals that explicit polynomial
  (Lemma~\ref{lem:linear-cr-220}).
\item \textbf{Not proved:} that the affine coset sum on the LHS of
  $(\dgg)$ equals the finite Cherednik--Ram sum. We articulated this as
  the ``coset-reduction gap'' (Section~\ref{sec:gap}) and provided a
  partial-evidence check that at least the diagram-automorphism $\pi$
  contribution collapses onto the finite sum (Section~\ref{ssec:pi-invariance}).
\end{itemize}

The single verifiable numerical statement of $(\dgg)_{3,3,(2,2,0)}$ ---
\emph{RHS is $(1+t) \cdot P_{(2,2,0)}$, which equals
$\sum_{S_3} T_u(x^{(0,2,2)})$} --- is therefore \emph{closed}. What
remains is the algebraic bridge from that finite sum to the literal
affine coset sum, which is the sharp form of $(\dgg)$.

Next-cycle probe: build the Bernstein $Y^\lambda$ operators from
the extended affine Hecke generators
$\wt T_w = T_u \cdot Y^\lambda$ for $\lambda \in Q/3Q$, and directly enumerate
$54$ representatives. This is the most promising path to actual verification
of $(\dgg)$ from the LHS side.

Companion scripts (see PROVE.md Step 5 ship section):
\verb|verify_linear_cr_220.py|, \verb|verify_pi_invariance.py|,
\verb|enumerate_affine_words.py|.

\begin{thebibliography}{9}

\bibitem{Cherednik}
I.~Cherednik, \emph{Double affine Hecke algebras}, LMS Lecture Note Series 319,
Cambridge Univ.\ Press, 2005.

\bibitem{Clio2026Boundary}
Clio (2026-07-21), memo
\emph{connections/2026-07-21-ter-cylindric-HL-interior-vs-boundary.md}.

\bibitem{ClioLinear}
Clio (2026-07-20), \emph{Demazure operators, Hall--Littlewood via Hecke,
and the transfer-op vs Demazure comparison},
\verb|~/projects/proofs/2026-07-20-demazure-basics-and-cherednik-ram-verification.pdf|.

\bibitem{ClioPieriAffine}
Clio (2026-07-24), \emph{The Pieri-level affine Cherednik--Ram theorem
is a corollary of the affine Cherednik--Ram conjecture via minuscule-character
centrality}, \verb|~/projects/proofs/2026-07-24-pieri-affine-CR-from-vdez.pdf|.

\bibitem{CR}
I.~Cherednik and A.~Ram, \emph{Double affine Hecke algebras and Macdonald's
conjectures}, Ann.\ Math.\ 141 (1995), 191--216.

\bibitem{Korff2013}
C.~Korff, \emph{Cylindric Hall--Littlewood polynomials and the affine flag
variety}, Comm.\ Math.\ Phys.\ 318 (2013), 173--246; arXiv:1110.6356.

\bibitem{Macdonald1995}
I.~G.~Macdonald, \emph{Symmetric Functions and Hall Polynomials},
2nd ed., Oxford Univ.\ Press, 1995.

\bibitem{Macdonald2003}
I.~G.~Macdonald, \emph{Affine Hecke Algebras and Orthogonal Polynomials},
Cambridge Tracts in Math.\ 157, Cambridge Univ.\ Press, 2003.

\bibitem{vDEZ2021}
J.F.~van~Diejen, E.~Emsiz, I.N.~Zurri\'an, \emph{Affine Pieri rule for
periodic Macdonald spherical functions and fusion rings}, Adv.\ Math.\ 392
(2021), 108027; arXiv:2305.01931.

\end{thebibliography}

\end{document}
